% Gerado por scripts/etapa3_comparacao.py — requer \usepackage{booktabs}
\begin{table}[htbp]
\centering
\small
\caption{Diferenças entre modelos, com intervalo de 95\% de confiança}
\label{tab:comp_diferencas}
\begin{tabular}{lrrrrrrr}
\toprule
Comparação & MAE & Pinball 80\% & Pinball 90\% & Pinball 95\% & Custo 1:4 & Custo 1:9 & Custo 1:19 \\
\midrule
XGBoost vs Holt-Winters & -1,3 [-1,6; -1,0] & -0,5 [-0,7; -0,2] & -0,2 [-0,5; +0,1] & +0,0 [-0,3; +0,3] & -0,5 [-0,8; -0,3] & -0,2 [-0,5; +0,1] & +0,0 [-0,3; +0,3] \\
Combinação vs Holt-Winters & -1,4 [-1,6; -1,2] & -0,9 [-1,1; -0,7] & -0,7 [-0,8; -0,5] & -0,5 [-0,7; -0,4] & -0,9 [-1,1; -0,7] & -0,6 [-0,8; -0,4] & -0,5 [-0,7; -0,3] \\
Combinação vs XGBoost & -0,1 [-0,2; +0,1] & -0,4 [-0,6; -0,3] & -0,5 [-0,6; -0,3] & -0,5 [-0,7; -0,4] & -0,4 [-0,5; -0,2] & -0,4 [-0,5; -0,2] & -0,5 [-0,7; -0,3] \\
\bottomrule
\end{tabular}
\par\smallskip
{\footnotesize Nota: teste de 90 dias (26/01 a 24/04/2016), 5.748 séries, quantis empíricos, estoque = quantil arredondado para o inteiro mais próximo. Diferença \% do primeiro modelo em relação ao segundo (negativo = primeiro melhor) e intervalo de 95\% de confiança entre colchetes, com o item como unidade (1.437 itens, cada um nas 4 lojas). Intervalo que não contém zero indica diferença distinguível do acaso.}
\end{table}
